\chapter{生命是一种怎样的化学系统}

\begin{kaichang}
抓一把干绿豆，铺在浸湿的纱布上，每天洒水。第三天，绿豆裂开，钻出白嫩的根；第五天，它变成一把挺立的豆芽。

现在请你认真回答一个问题：那把干绿豆，和这把豆芽，哪一个更“混乱”？

直觉可能会说：干绿豆乱糟糟堆在一起，豆芽却根朝下、芽朝上，排列整齐，还长出了精巧的结构——豆芽明显更“有秩序”。可我们在第 28 章刚刚认识了一条铁律：孤立系统里，一切过程都让熵增加，万物只会走向更混乱，绝不自发走向更有序。一颗泡在水里的豆子，没有设计师，没有图纸，凭什么越长越精巧？

它是自然规律的例外吗？先写下你的猜想。这一章结束时，我们回来对答案——而这一章之后，你就正式从“化学”走进了“生命”。
\end{kaichang}

\section{先看几件摆在眼前的事}

在给出任何理论之前，先把眼睛能看、手能摸、温度计能测的事实摆出来。拿生命和死物对比，差别一目了然：

\begin{itemize}[itemsep=2pt]
\item 一把豆芽每天都在长高、变重，结构一天比一天复杂；一颗玻璃珠放在同样的盘子里，一万年也不会变。
\item 一碗米饭放两天就发馊——有看不见的东西在“吃”它；而真空包装、彻底灭菌的罐头可以放好几年。
\item 你安静坐着，体温也稳稳维持在 $37\,^{\circ}\mathrm{C}$ 上下，冬天夏天差不多；一块石头只会随波逐流地变成环境温度。
\item 你的手划破了，过几天伤口自己长好；毛衣划破了，只会越破越大。
\item 秋天落叶归根，几年后就“消失”在土壤里——又被看不见的东西拆成了原料。
\end{itemize}

这些现象还可以测量：把一只小动物（或者一颗正在发芽的种子）放进密闭容器，氧气会减少，二氧化碳会增加，温度计会微微上升。活着的东西在持续地“吃进”什么、“排出”什么，还持续地发热。这是生命最扎眼的特征：\textbf{它不是一件东西，而是一股持续不停的流}。

\begin{shiyan}
\textbf{发一盘豆芽，看生命“开工”}

材料：干绿豆一小把、一个浅盘或沥水篮、干净的纱布或厨房纸、清水、一支温度计（可选）、一台厨房秤（可选）。

步骤：
\begin{enumerate}[itemsep=1pt]
\item 取两份等量的干绿豆（比如各 \ud{20}{g}），一份煮熟放凉作为对照，一份保持生的。
\item 分别铺在两个铺了湿纱布的浅盘里，标注“生”“熟”，放在室温避光处。
\item 每天同一时间洒水、观察，连续 5～7 天。记录：豆粒是否膨胀、是否破皮出根、纱布是否有异味；有秤就称一称湿重的变化；有温度计就在豆子堆里埋几分钟，比较两盘的温度。
\end{enumerate}

观察点：生绿豆吸水膨胀、破皮、生根、长芽，几天后盘上会出现微微的温热（发芽旺盛时）；煮熟的那份只会慢慢发黏、长霉。想一想：霉又是谁？它也在“开工”——只不过开的是另一套工。

安全提示：这是安全的家庭观察。注意勤洗手；发霉的豆子不要尝，直接扔掉并清洗容器。
\end{shiyan}

\section{把镜头推进到细胞}

豆芽为什么会动、会长、会发热？把镜头往里推：豆芽的每一小段，都由亿万个微小的“房间”组成——细胞（第 57 章会带你逐个房间参观）。每个细胞被一层薄薄的膜包着（第 51 章专门讲这层膜），膜内不是一潭静止的汤，而是一座昼夜开工的工厂：营养分子从膜外运进来，被拆小、被改造、被重新拼装，废物和热量被排出去。

规模有多夸张？一个普通细胞里，每秒钟发生的化学反应以百万计；你全身的细胞每秒钟合成的蛋白质分子数以万亿计。而这一切不是乱哄哄的混战：什么分子在什么时间、什么位置、与谁反应，受到精确的控制——负责调度这些反应的，是第 49 章要讲的酶。

回头看贯穿全书的五个问题，现在问的是“生命”这个对象：

\begin{enumerate}[itemsep=2pt]
\item 它由什么组成？——水、蛋白质、糖、脂质、核酸，归根结底是碳、氢、氧、氮、磷、硫等原子。
\item 怎样连接排列？——共价键搭成分子，分子再组装成细胞结构，层次井然。
\item 什么规则决定它能否变化？——和烧煤、生锈一样，是能量与熵的规则，没有例外。
\item 变化多快、怎样受控？——靠酶和调节网络，快慢收放自如。
\item 它怎样被保存、复制、调节和选择？——这就是“活着”的核心，本章的主线。
\end{enumerate}

注意第三个问题的答案：我们没有给生命另立一套物理定律。豆芽和玻璃珠服从同样的规则。那么，豆芽凭什么“逆着混乱生长”？答案藏在那条铁律的一个限定词里。

\section{那条铁律只管“孤立系统”}

第 28 章说过：\textbf{孤立}系统——既不与外界交换物质、也不交换能量的系统——的熵只增不减。关键词是“孤立”。

豆芽是孤立系统吗？显然不是：它吸水、吸氧气、从豆子里调取储存的营养，还向外排水汽、二氧化碳和热量。你也不是孤立系统：你吃饭、呼吸、排汗、散热。孤立系统在自然界几乎不存在，只是一个思想工具。真正的第二定律用在开放系统上，说法是：\textbf{系统加环境，总的熵只增不减}。至于系统内部？只要它向外排出的熵足够多，内部完全可以让熵减少——长出精巧的结构。

你家厨房里就有一台天天演示这件事的机器：冰箱。一杯水放进冷冻室，水分子排成整齐的冰晶——熵大大降低了，这“违反”第二定律吗？一点也不。摸摸冰箱背后或两侧的散热板：热的。冰箱压缩机把热量从内部搬到外部，排给厨房空气；把冰箱和房间算总账，熵还是增加的。冰箱内部那点“秩序”，是用外部更多的“混乱”换来的。

生命干的是同一件事，只是精巧得多。细胞是一座开放的工厂：吃进低熵的“原料”（结构复杂、能量集中的食物分子和氧气），排出高熵的“废料”（小而散的水、二氧化碳分子，以及四散的热运动）。\textbf{细胞内部维持低熵结构的代价，是让环境的熵增加得更多。}总账永远平衡，第二定律纹丝不动。

\begin{qiaoliang}
来算一笔真实的账。一个成年人一天大约消耗 2000 千卡的食物能量，这些能量最终几乎都以热的形式排给环境（你干体力活搬起的重物，最后也变成热）。换算一下：
\[2000\ \mathrm{kcal} \approx \ud{8.4\times 10^{6}}{J}\]
假设你的皮肤和周围环境大约是 $27\,^{\circ}\mathrm{C}$，即 $300\ \mathrm{K}$，这些热量带给环境的熵增加约为
\[\Delta S \approx \frac{Q}{T} = \frac{8.4\times 10^{6}\ \mathrm{J}}{300\ \mathrm{K}} \approx \ud{2.8\times 10^{4}}{J\cdot K^{-1}}\]
这个数字太抽象？找个参照：让 $0\,^{\circ}\mathrm{C}$ 的冰融化成 $0\,^{\circ}\mathrm{C}$ 的水，每千克大约吸热 \ud{3.3\times 10^{5}}{J}，熵增加约 \ud{1.2\times 10^{3}}{J\cdot K^{-1}}。也就是说，\textbf{你一天向环境排出的熵，相当于把 20 多千克冰融化成水所增加的熵}。你身体里长出新的细胞、修复组织、维持精巧结构所“省下”的熵，跟这笔排出的账比起来只是零头——总账上熵依然大幅增加。
\end{qiaoliang}

把视野再放大到整个地球。生物圈这台大机器的能量来源是太阳：太阳送来的可见光，每个光子能量高、数量相对少——这是低熵的能量流；地球把同样的能量以红外辐射还给太空，每个红外光子能量低、数量多得多（同样一份能量，摊给了大约二十倍数量的光子）——这是高熵的能量流。能量守恒，一进一出；熵呢？出去的热闹远超进来的。整个生物圈，包括每一根豆芽、每一片森林和你，都“挂”在这条从太阳流向太空的能量河上，靠它冲刷维持自己的秩序。

\begin{center}
\begin{tikzpicture}[scale=0.95]
\draw (0,0) rectangle (4.6,3);
\node at (2.3,1.5) {细胞（维持低熵结构）};
\draw[->,thick] (-3,2.2) -- (0,2.2);
\node at (-1.6,2.65) {低熵输入};
\node at (-1.6,1.85) {食物、氧气};
\draw[->,thick] (4.6,2.2) -- (7.8,2.2);
\node at (6.2,2.65) {高熵输出};
\node at (6.2,1.85) {二氧化碳、水、废物};
\draw[->,thick] (4.6,0.7) -- (7.8,0.7);
\node at (6.2,0.3) {热（最“廉价”的熵）};
\end{tikzpicture}
\end{center}

（这是人为的表示法：箭头画的是物质与能量的账目流向，不是细胞的真实长相。）

\section{生命递出的五张名片}

到这里，“生命是一台不违反第二定律的开放化学系统”已经能解释开场的现象了。但物理学家会追问：那么，一股瀑布、一团火焰也是开放的能量流，它们算生命吗？显然不算。区别在哪？生物学家清点下来，生命通常递出五张名片——

\begin{itemize}[itemsep=2pt]
\item \textbf{边界}。生命把自己和环境分开：细胞有一层膜，你有皮肤。边界之内维持着和外界不一样的成分与秩序，边界有选择地放行物质。火焰没有边界，风一吹就散。第 51 章会细讲这层膜怎样“把关”。
\item \textbf{代谢}。生命持续地进行化学反应：拆（分解营养、释放能量）也建（合成自身材料、储存能量）。水晶也会“长”，但它只是从溶液里捞出现成的分子堆上去，自己内部什么都不发生。第 52～55 章专门讲代谢这张网。
\item \textbf{复制}。生命能照着一份“说明书”复制自己：一颗绿豆长成豆苗，豆苗结出新绿豆。火焰也能“蔓延”，但每次蔓延都不像谁，没有一份被抄写的东西。这份说明书——DNA——是第三册的主角。
\item \textbf{变异}。抄写偶尔会抄错，环境也会改动说明书，于是后代和亲代有差别。没有变异，下一条就无从谈起。
\item \textbf{演化}。在一代又一代的筛选下，有利于生存繁殖的变异被留下来，种群随时间改变——这是第三册后半部的主线。火焰和瀑布都没有“代”，更谈不上演化。
\end{itemize}

五张名片凑齐，才是我们熟悉的那种“活”。注意这套判据是\textbf{描述性的清单，不是自然刻在石头上的定义}——这一点，等会儿讨论病毒时有大用。

\section{这台化学机器用什么原料}

说生命是化学系统，就要交代它的原料单。把你身体里的原子过一遍秤：氧约占六成多，碳约两成，氢约一成，氮约百分之三，磷、硫和其余几十种元素加起来只占百分之几。但“少”不等于“不重要”——螺丝钉也比车身轻得多。

\begin{itemize}[itemsep=2pt]
\item \textbf{水}不是元素，却是戏份最重的角色：你体重的六成以上是水。它是细胞的溶剂、运输介质、体温缓冲垫，还亲自参加许多反应。第 23 章讲过它那些反常的脾气——能溶解极性物质、比热大、结冰膨胀——没有一条是白长的，每一条都被生命用上了。
\item \textbf{碳}是所有生物大分子的骨架。第 36 章讲过原因：碳原子恰好能形成四个共价键，能连成链、环和网，稳定又多变。糖、脂质、蛋白质、核酸，全是碳搭的。
\item \textbf{氮}藏在蛋白质和核酸的分子里。空气里 78\% 是氮气，你却一口都用不上——\chem{N_2} 太稳定，得靠细菌和工业把它“拆开”变成可用的形态，这是第 14 章埋过的伏笔，后面讲物质循环还会回来。
\item \textbf{磷}是两件关键装备的骨架：一是遗传物质 DNA 的“链条龙骨”，二是能量通货 ATP（第 50 章会讲，顺便澄清“高能键”这个流行说法错在哪）。你的骨骼和牙齿也靠磷酸钙撑着。
\item \textbf{硫}戏份少但不可替代：某些氨基酸里有它，两个硫原子能拉成“二硫键”，像订书钉一样把蛋白质的长链订成固定形状（第 48 章）。你烫头发的原理，就是先拆再订这些硫“订书钉”。
\item 剩下的“微量元素”——铁、锌、碘、镁……——总量微不足道，却各管一个关键岗位：血红蛋白中心的铁抓住氧气，叶绿素中心的镁抓住阳光（第 15 章）。
\end{itemize}

\begin{table}[htbp]
\centering
\begin{tabular}{lll}
\toprule
角色 & 在生命里的主要差事 & 哪里细讲 \\
\midrule
水 & 溶剂、运输、控温、参与反应 & 第 23 章 \\
碳 & 一切生物大分子的骨架 & 第 36 章 \\
氮 & 蛋白质和核酸的核心成分 & 第 48 章 \\
磷 & 遗传物质与能量通货的骨架 & 第 50 章 \\
硫 & 帮蛋白质固定形状的“订书钉” & 第 48 章 \\
\bottomrule
\end{tabular}
\caption{生命的几位主力角色。这张表只是索引，后面每一章展开其中一位。}
\end{table}

看这张单子，你大概已经闻到了第一篇到第五篇的味道：生命没有发明新元素，也没有发明新键，它只是在周期表的一小角里，用你已经认识的规则，搭出了会自我维护的结构。

\section{把这笔账写成式子}

现在轮到符号登场了。细胞“吃进低熵、排出高熵”这笔总账，可以写成你已经熟悉的语言——化学方程式。以葡萄糖为例，细胞呼吸的总账是：
\[\chem{C_6H_{12}O_6} + \chem{6O_2} \rightarrow \chem{6CO_2} + \chem{6H_2O}\]

左边：一个结构复杂的糖分子，加六个氧分子；右边：十二个小而简单的分子，外加释放的能量。注意看清方向：断键吸能、成键放能（第 27 章），这步总账放能，是因为生成物里新形成的键比反应物里的“更划算”。一个葡萄糖分子是集中、有序的；把它拆成一群小分子、让能量摊成热，就是熵增加的康庄大道——这正是细胞维持自身秩序时，向外支付的那笔“熵账单”。

但要马上泼两盆冷水。第一，这只是\textbf{总账}：细胞绝不让葡萄糖“一把烧掉”，而是拆成几十个受控的小步骤，一步步把能量“兑零”存进 ATP ——第 52、53 章会带你走完这条流水线。第二，植物的光合作用乍看像把这个式子倒过来写，但两者\textbf{并不是简单的互逆过程}，而且释放的氧气来自水，不是来自二氧化碳——这个常见的误会，第 54 章要专门拆掉。

\begin{zhengju}
\textbf{豚鼠和炭火：生命服从化学规律的第一次称量}

18 世纪末，化学还流行一种看法：生物体内有某种“生命力”，生物的变化不受无机界规律约束。是就是，不是就不是，得拿实验称一称。1780 年代，法国化学家拉瓦锡和数学家拉普拉斯造了一台巧妙的仪器——冰量热计：一个夹层容器，夹层里填满碎冰，内腔放待测对象；对象放出多少热，就融化多少冰，称一称滴出来的水，热量就有了。

他们先放进一只豚鼠：十个小时，它呼吸、产热，融化了不少冰，同时让内腔里的空气多了二氧化碳。然后换成烧炭：精确称量，让炭燃烧产生\textbf{同样多}的二氧化碳，结果融化的冰几乎一样多。结论大胆而朴素：\textbf{呼吸就是一种缓慢的燃烧}——同样消耗氧、同样产生二氧化碳、同样放热，只是慢得多、温和得多。维持体温的热量，不来自什么神秘的生命力，而来自可以用方程式记账的化学反应。

这个故事也有该修正的部分，这让它更可信。拉瓦锡以为“燃烧”主要发生在肺部；几十年后人们才搞清楚，氧化发生在全身的细胞和组织里，肺只是换气站。而他用的“热质说”框架，也要等到 19 世纪 40 年代能量守恒定律确立后才被彻底纠正——热不是一种流来流去的物质，而是能量的一种形态。科学就是这样：结论可以先对，解释慢慢补；能被后来的证据修正，恰恰说明它走在对的路上。

顺带一提，拉瓦锡的这台冰量热计，就是你家里那只动物（和你自己）遵守化学规律的第一张收据。此后两百年，从肌肉收缩到神经放电，生命活动一项一项被换算成了能量和物质的账目——没有一项需要“例外条款”。
\end{zhengju}

\section{这个模型的边界，和病毒这个“刺头”}

“开放系统吃进低熵、排出熵”这个模型，什么时候好用？解释你为什么会饿、会热、要呼吸，解释豆芽为什么能长，解释伤口为什么能愈合——都好用。它把生命安安稳稳地放进了物理定律的版图。

但它也有边界，边界处蹲着一个著名的刺头：\textbf{病毒}。

拿五张名片去核对病毒：它有边界吗？有一层蛋白质外壳，有时还偷来一片膜——勉强算有。它会变异、会演化吗？太会了，流感病毒年年变脸。它能复制吗？能——但全靠劫持宿主细胞的工厂，自己一个零件都造不了。最关键的一张：\textbf{代谢}。病毒不吃、不喝、不呼吸、不排熵，离开宿主就是一颗静止的分子颗粒，不维持任何低熵的内部秩序，可以冻在冰箱里几年“死而不僵”。

所以病毒算不算生命？诚实的回答是：\textbf{取决于你把名片清单的及格线画在哪}。把“自主代谢”设为及格线，病毒出局；把“能复制会演化”设为及格线，病毒入选。这不是和稀泥——它说明“生命”这个词是人类画的圈，自然界只负责提供连续的渐变：从病毒、类病毒（一小段核酸）、朊病毒（一个会传染的错误折叠蛋白质），到细菌、到你，界限本来就模糊。类似的还有休眠的种子和冻存的细胞：代谢低到几乎测不出，它们“活着”吗？靠名片清单，你可以给出言之成理的答案；指望自然界给你一条刀切般的分界线，会失望的。

这就是本章模型的使用说明：它告诉你生命\textbf{怎样}运行，不替你裁决“什么算生命”。前者是科学问题，后者一半是定义问题。

\begin{wuqu}
“第二定律说万物走向混乱，生命却越来越有序，所以生命（和进化）违反自然规律，必须有超自然的解释。”

反例就在你家厨房和窗台上：冰箱里的水冻成整齐的冰晶，冬天窗玻璃上长出精巧的冰花——它们都比原来“有序”，没人觉得需要奇迹，因为冰箱向外排热、窗玻璃向冷空气排热，总账上熵照增不误。生命同理：它是\textbf{开放系统}，吃进低熵的食物和阳光，排出高熵的废物和热，整个“生物加环境”的总熵一直在增加。第二定律从诞生那天起就只管辖孤立系统，拿它去审判豆芽，是用错了说明书。顺便说，连“进化让生物越来越复杂”本身都不是定律的要求——进化没有预定方向，复杂化只是局部故事里常见的情节，不是宇宙的目标。
\end{wuqu}

\section{再看那把豆芽}

回到开场的问题：干绿豆和豆芽，哪个更混乱？

现在可以给出一个有分寸的答案了。豆芽内部确实比一把干豆子更“有序”：细胞分化、组织成形、结构精巧——局部的熵降低了。但这幅画面少了一半：豆芽每时每刻在吸水、消耗豆子里储存的营养、吸入氧气，排出二氧化碳、水汽和热量。把“豆芽加它周围的环境”作为一个整体来记账，熵每时每刻都在增加。那把看似“逆着混乱”生长的豆芽，其实是一台把储存的化学秩序兑换成自身结构、同时向环境大排熵的机器——和冰箱同一原理，只是精致了亿倍。

顺带解开实验里的小悬念：煮熟的绿豆为什么不发芽，只会发霉？加热破坏了它细胞里蛋白质机器的形状（第 48 章叫“变性”），代谢的工厂停了工，名片一张也递不出来——它从开放系统退化成了一堆普通有机物，于是轮到另一批活着的开放系统（霉菌）来拆解它。生与死的区别，在化学上看，就是\textbf{那股自我维持的流还在不在}。

\section{继续深挖}

这根“能量流”的线索可以一直挖下去。医学里，发高烧、甲状腺功能亢进，本质都是代谢这团火的“阀门”出了问题；临床上测量基础代谢率，就是在给人这台机器读功率表。运动科学里，你跑步时浑身发烫、大汗淋漓，正是总账式子里那笔“热量支出”在体感上的样子。生态学里，整条食物链可以看成一条逐级传递的能量河：草抓住阳光，羊吃草，狼吃羊，每一级都有大笔能量以热的形式“漏”给环境——这就是为什么食物链通常只有四五个环节，顶端的掠食者总是稀少。甚至生命起源问题（第 56 章）也绕不开本章的框架：最早的生命雏形，很可能就诞生在海底热泉口那种天然的能量流里——先有一股可搭便车的流，才谈得上学会自己开流。

而眼下，最顺理成章的下一个问题是：细胞吃进的食物分子，具体长什么样？凭什么糖能当燃料、当建材、还能当“身份标签”？翻开下一章，我们从最常见的那个分子——葡萄糖——讲起。

\begin{tiaozhan}
想再深一层？第 28 章提过，熵在微观上的定义是玻尔兹曼公式 $S = k\ln W$，其中 $W$ 是系统微观排法的数目。用这个眼光看本章：一个活细胞把原子组织成高度特定的排布，$W$ 很小，熵很低；它每合成一个这样的结构，都必须通过代谢把更多能量摊散给环境，让环境的 $W$ 增加得更多。1944 年，物理学家薛定谔在小册子《生命是什么》里把这个想法讲给世界听，他造了个词叫“负熵”——说生命“以负熵为食”。这个词很抓人，但容易误导：熵不是可以被吃掉的东西，今天更准确的表述是，生命摄取的是\textbf{自由能}（第 28 章），并把熵排出体外。这本小册子影响了后来一批转向生物学的物理学家——包括 DNA 双螺旋故事里的关键人物，那是第三册要讲的故事。跳过本框不影响主线。
\end{tiaozhan}

\section*{练一练}

\begin{sikaoti}
\item 画一幅“细胞物质流图”：方框代表细胞，箭头标出进入和离开细胞的物质与能量，并在每条箭头旁注明它是低熵还是高熵。用一句话说明：为什么这张图不违反热力学第二定律？
\item 冰箱冷冻室里的水结成整齐的冰晶（熵降低）。这违反第二定律吗？请写出“把谁和谁算总账”，并解释冰箱背部为什么发热。
\item 一碗米饭大约提供 500 千卡能量。假设这些能量最终全部以热的形式排给 $300\ \mathrm{K}$ 的环境，估算环境熵增加多少（用 $\Delta S \approx Q/T$，$1\ \mathrm{kcal}\approx\ud{4.2\times 10^{3}}{J}$）。再换算成“相当于融化多少千克 $0\,^{\circ}\mathrm{C}$ 的冰”。
\item 用本章的“五张名片”逐一核对：一团篝火、一颗休眠三年的种子、一个流感病毒、一粒冷冻保存的细菌。哪些项目它们各自拿不到？你的“及格线”画在哪里，为什么？
\item 干绿豆可以存放几年不变，豆芽放两三天就腐烂。用“开放系统、水、代谢”三个关键词，解释这个差别。
\item 本章说“生命没有发明新元素、新键，只是用旧规则搭出了会自我维护的结构”。从第一册任选一条规则（如成键放能、平衡移动、催化剂），举一个生命正在使用它的例子。
\end{sikaoti}
